\scanpage{57}
Le groupe quotient
\[
\operatorname{Aut}(I,I')/
\underbrace{\operatorname{Aut}^{\unsure{0+|o+}}(I,I')}_{
\substack{\text{autom. qui}\\\text{respectent }I,I'\\\text{et l'orientation}\\
\text{de }I,\ \text{i.e. de }I'}}
\]
est isomorphe à \uline{$\mathbb Z/4\mathbb Z$} sauf erreur.
Mais si on \deleted{\unsure{est|veut}} se donne une
\deleted{\unsure{comme}} \deleted{\unsure{str} biicosaèdre orienté}
orientation de $I$ et une de $I'$ (pour
définir une notion de biicosaèdre \uline{orienté})
\deleted{on \unsure{trouve}}\inserted{d'où} $\unsure{\vec I|\overline I}$ et $\vec I'$, cela nous permet
de choisir \unsure{l'un|\deleted{un}} des icos. de la structure
biicosaédrale de façon canonique -- \unsure{nommément|comme}
celui pour lequel l'orientation de l'autre
soit ``associée''. Donc on ne trouve pas
la bonne notion !

Par contre, il y a un groupe $\mathbb Z/4\mathbb Z$
qui opère sur la catégorie des
icosaèdres orientés, dont le générateur est
\deleted{le passage de} est le passage de $\vec I$ à $\vec I^\circ$.
Le carré du générateur est le passage
de $I$ à $-I$ -- correspond à l'ext.
triviale de $\mathbb Z/2\mathbb Z\subset\mathbb Z/4\mathbb Z$ par
$\operatorname{Aut}(\vec I)^+$, savoir le groupe des autom
de \deleted{Aut} $I$ (ne respectant pas l'orientation)
% archive sheet number: 56
